% 高校生ガイド v2
\documentclass[12pt,a4paper]{article}

\usepackage{fontspec}
\usepackage{xeCJK}
\setCJKmainfont{Hiragino Mincho ProN}
\setCJKsansfont{Hiragino Kaku Gothic ProN}
\usepackage[margin=2.5cm]{geometry}
\usepackage{amsmath,amssymb}
\usepackage{hyperref}
\usepackage{tcolorbox}
\usepackage{booktabs}
\usepackage{fancyhdr}

\tcbuselibrary{breakable,skins}

\newtcolorbox{storybox}[1][]{colback=blue!5,colframe=blue!60!black,
  fonttitle=\bfseries\sffamily,title={#1},breakable,arc=3mm}
\newtcolorbox{mathbox}[1][]{colback=green!5,colframe=green!50!black,
  fonttitle=\bfseries\sffamily,title={#1},breakable,arc=3mm}
\newtcolorbox{resultbox}[1][]{colback=yellow!8,colframe=orange!80!black,
  fonttitle=\bfseries\sffamily,title={#1},breakable,arc=3mm}
\newtcolorbox{honestbox}[1][]{colback=gray!8,colframe=gray!50!black,
  fonttitle=\bfseries\sffamily,title={#1},breakable,arc=3mm}

\setlength{\parskip}{0.8em}
\setlength{\parindent}{0pt}

\pagestyle{fancy}
\fancyhf{}
\rhead{Wright Brothers}
\lhead{高校生ガイド v2}
\cfoot{\thepage}

\begin{document}

\begin{center}
{\Huge\bfseries 負のエネルギーの壁を}\\[0.2cm]
{\Huge\bfseries 位相で突破する}\\[0.8cm]
{\Large ギターとビール瓶、超伝導、そしてワープ}\\[1cm]
{\large 高校生のためのガイド}\\[0.5cm]
{\large Wright Brothers Research Program, 2026年4月}
\end{center}

\vspace{0.5cm}
\tableofcontents
\newpage

%====================================================================
\part{3つの事実}
%====================================================================

\section{事実 1: 何もない空間に力がある}

\begin{storybox}[カシミール効果]
2枚の金属板をナノメートルの間隔で並べると、
何もないはずの板の間に\textbf{引力}が生じる。

板の間の真空は外の真空より\textbf{エネルギーが低い}。
エネルギーが低い → 板が引き寄せられる。

この力の大きさは:
$$F = -\frac{\pi^2 \hbar c}{240 \, a^4} \times A$$

ここで $a$ = 板の間隔、$A$ = 板の面積。

この「240」はリーマンゼータ関数の値 $\zeta(-3) = 1/120$ から来る。
（$240 = 2/\zeta(-3)$、2 は電磁波の偏光方向の数。）

\textbf{1997年に実験で確認。理論との一致 1\% 以内。}
\end{storybox}

\begin{resultbox}[意味]
「$\zeta(-3) = 1/120$」は計算のトリックではない。
\textbf{物理的に測定可能な量}。板にかかる力がそう言っている。
\end{resultbox}

\section{事実 2: 符号は変えられる}

\begin{storybox}[ギターとビール瓶]
ギター（両端固定）: 全倍音 → 真空エネルギー $\propto \zeta(-1) = -1/12$（\textbf{負}）

ビール瓶（片端開放）: 奇数倍音のみ → 真空エネルギー $\propto +1/12$（\textbf{正}）

片端を開けるだけで符号が逆になる。

2009年（Nature）と 2017年（Nature Photonics）に
カシミール\textbf{斥力}が実験的に実証された。
引力（負のエネルギー）が斥力（正のエネルギー）に変わる。
\end{storybox}

\begin{resultbox}[意味]
真空エネルギーの符号は\textbf{境界条件で制御できる}。
これは確立された物理学。
\end{resultbox}

\section{事実 3: 宇宙の加速膨張}

\begin{storybox}[暗黒エネルギー]
宇宙は加速しながら膨張している。
この加速の原因は「暗黒エネルギー」と呼ばれ、
真空そのもののエネルギー密度 $\rho_\Lambda \approx 6 \times 10^{-10}$ J/m³ で記述される。

この値は $\rho_\Lambda \propto \zeta(-3)$ と整合的
（カシミール力と同じ $\zeta(-3)$ が宇宙論にも現れる）。
\end{storybox}

\section{自然な推論: 第4ステップ}

\begin{resultbox}[3つの事実からの論理的帰結]
\begin{enumerate}
\item $\zeta(-3)$ がカシミール力を決める → 実験確認済み
\item $\zeta \to \zeta_{\neg 2}$ で符号が反転する → 実験実証済み
\item 暗黒エネルギーも $\zeta(-3)$ に従う → 観測と整合的
\end{enumerate}

$\therefore$ 暗黒エネルギーの符号も $\zeta \to \zeta_{\neg 2}$ で反転可能。

\textbf{これを禁じる確立された物理法則はない。}
\end{resultbox}

%====================================================================
\part{壁: なぜ簡単にはいかないか}
%====================================================================

\section{表面効果の限界}

\begin{mathbox}[カシミールエネルギー密度]
板の間のエネルギー密度:
$$\rho = -\frac{\pi^2 \hbar c}{720 \, a^4}$$

板を離すと密度は急激に下がる:
\begin{center}
\begin{tabular}{ll}
$a = 1$ nm: & $\rho \sim 10^{12}$ J/m³（巨大） \\
$a = 1\,\mu$m: & $\rho \sim 10^{-4}$ J/m³（微小） \\
$a = 1$ m: & $\rho \sim 10^{-24}$ J/m³（ゼロ同然） \\
\end{tabular}
\end{center}

\textbf{板を離すと効果が消える。} ナノの世界でしか大きくない。
\end{mathbox}

\section{Ford-Roman の壁}

\begin{mathbox}[量子エネルギー不等式]
Ford と Roman（1995, 1997）は次を証明した:

\textbf{「負のエネルギーの量 $\times$ 持続時間 $\lesssim \hbar$」}

つまり: 大きな負のエネルギーは一瞬しか持てない。
長時間持てる負のエネルギーは極めて微小。

ワープに必要: 巨大な負のエネルギーを持続的に。
Ford-Roman: それは不可能。

\textbf{これが「70桁の壁」。} カシミールの $10^{-8}$ J と
ワープの $10^{62}$ J の間の絶望的な差。
\end{mathbox}

%====================================================================
\part{突破口: 位相的バイパス}
%====================================================================

\section{Ford-Roman の仮定}

\begin{resultbox}[見落とされた仮定]
Ford-Roman が仮定しているのは:
\textbf{「通常の真空から連続的に変形して到達できる状態」}

輪ゴムに例えると:
Ford-Roman は「ゴムを引っ張る」操作しか考えていない。
「ゴムを切って結び直す」（位相的操作）は考えていない。
\end{resultbox}

\section{超伝導の教訓}

\begin{storybox}[BCS 超伝導]
1900年代の物理学者は「電気抵抗をゼロにはできない」と思っていた。

Drude 理論: 温度を下げると抵抗は減るが、決してゼロにはならない。
不純物があるから。

BCS 理論（1957年）: ある温度で\textbf{相転移}が起き、
電子がペアを組み（クーパー対）、全く新しい量子状態に遷移する。
この状態では抵抗は正確にゼロ。

Drude は間違っていなかった。「連続的に冷やす」限り、
抵抗は確かにゼロにならない。
しかし\textbf{相転移}（不連続なジャンプ）が起きた瞬間、
Drude の前提が崩れた。
\end{storybox}

\begin{resultbox}[同じ構造]
Ford-Roman は間違っていないかもしれない。
「連続的に変形する」限り、大きな負のエネルギーは確かに無理。

しかし\textbf{位相的転移}（K理論の巻き付き数のジャンプ）が起きたら？

$$K_1(\mathbb{Z}) = \mathbb{Z}/2 \;\xrightarrow{\text{ジャンプ}}\; K_1(\mathbb{Z}[1/2]) = \mathbb{Z}/2 \times \mathbb{Z}$$

超伝導の抵抗ゼロが Drude の外にあったように、
位相的に保護された負のエネルギーは Ford-Roman の外にある。
\end{resultbox}

\section{体積効果への格上げ}

\begin{storybox}[壁のペンキから部屋の水へ]
カシミール効果 = 壁のペンキ塗り（表面効果）。
板を離すとペンキの面積が増えるが、密度は $1/a^4$ で薄まる。

もし位相的転移が起きると = 部屋を水で満たす（体積効果）。
密度は板間隔によらず一定。

\begin{center}
\begin{tabular}{lcc}
\toprule
& 表面効果 & 体積効果 \\
\midrule
密度 & $\rho \propto 1/a^4$ & $\rho =$ 一定 \\
$a = 1$ m & $10^{-24}$ J/m³ & $10^{12}$ J/m³ \\
体積 1 m³ でのエネルギー & $10^{-24}$ J & $10^{12}$ J \\
\bottomrule
\end{tabular}
\end{center}

\textbf{$10^{36}$ 倍の差。}
\end{storybox}

%====================================================================
\part{たった一つの実験}
%====================================================================

\section{何を測るか}

\begin{resultbox}[最もシンプルな実験]
ビール瓶型（DN 境界条件）のカシミール力を、
板間隔 $a$ を変えながら精密に測定する。

力のべき乗則をフィット: $F \propto a^{-\alpha}$

標準 QED: 全ての $a$ で $\alpha = 5$（一直線）

位相的予測: ある $a_c$ で $\alpha$ が 5 から折れ曲がる

$a_c$ がどこかは予測しない。
予測するのは\textbf{「折れ曲がりが存在する」}こと。

\begin{center}
\begin{tabular}{ll}
折れが\textbf{ない} → & 位相的保護は（測定範囲で）存在しない \\
折れが\textbf{ある} → & \textbf{真空の位相的転移の発見} \\
\end{tabular}
\end{center}
\end{resultbox}

\section{誰がやるか}

\begin{storybox}[Princeton 大学]
2017年、Princeton大学のチームが Si-MEMS デバイスで
カシミール斥力を実測した（Nature Photonics）。

同じ装置で、板間隔 $a$ を系統的に変えて力のべき乗を測れば、
$\alpha = 5$ からの逸脱の有無が判定できる。

\textbf{新しい装置の発明は不要。既にある装置で追加データを取るだけ。}
\end{storybox}

%====================================================================
\part{もし位相的転移が見つかったら}
%====================================================================

\begin{resultbox}[4つの帰結]
\begin{enumerate}
\item \textbf{制御可能なダークエネルギー}:
境界条件を変えるだけで真空エネルギーの符号をスイッチできる。

\item \textbf{巨視的な負のエネルギー}:
位相的に保護された定数密度 × 体積 = 発電機1台で巨大なエネルギー。

\item \textbf{Ford-Roman の壁の突破}:
持続的な負のエネルギーが位相的相で実現。70桁のギャップが消える。

\item \textbf{ワープドライブへの道}:
十分な負のエネルギーがあれば、Alcubierre のワープ計量が
物理的に実現可能に。
\end{enumerate}
\end{resultbox}

%====================================================================
\part{正直な話}
%====================================================================

\begin{honestbox}[確率の話]
\textbf{折れ曲がりが見つかる確率}: 低い。

標準 QED は70年以上、全ての精密テストに合格してきた。
$1/a^4$ のスケーリングが破れる兆候は今のところゼロ。

しかし:
\begin{itemize}
\item DN 配置でのスケーリングは\textbf{一度も精密測定されていない}
\item 位相的転移は「連続的な測定」では見えない種類の効果
\item BCS 超伝導も発見前は「ありえない」と思われていた
\item 測定して「ない」とわかることにも科学的価値がある
\end{itemize}

\textbf{もし見つかったら}: 物理学の歴史が変わる。\\
\textbf{もし見つからなければ}: DN カシミールのスケーリング則を初めて精密測定した論文として残る。

\textbf{どちらに転んでも科学は前に進む。}
\end{honestbox}

\vfill

\begin{center}
\rule{10cm}{0.4pt}\\[0.5cm]
{\Large\bfseries 壁があるなら、越えるのではなく、}\\[0.2cm]
{\Large\bfseries 位相を変えて壁ごと消す。}\\[0.3cm]
{\small Wright Brothers Research Program, 2026}
\end{center}

\end{document}
